%======================================================================
\section{Firmware}\label{sec:firmware}
%======================================================================
\subsection{Ambiente e organização}

O firmware das duas placas é escrito em C com o ESP-IDF 6.0.1~\cite{espidf}, que já traz o
FreeRTOS, as pilhas de Wi-Fi e TCP/IP, o cliente MQTT e os drivers dos periféricos. São dois
projetos, um por placa, e um conjunto de componentes comuns que as duas compilam
(Tabela~\ref{tab:componentes}). Cada placa passa os seus pinos e tópicos aos componentes por
uma estrutura de configuração; nada é copiado de uma placa para a outra.

\begin{table}[H]
\centering
\caption{Componentes do firmware.}
\label{tab:componentes}
\small
\begin{tabular}{@{}l>{\raggedright\arraybackslash}p{11.6cm}@{}}
\toprule
Componente & Responsabilidade \\
\midrule
\multicolumn{2}{@{}l}{\textit{Comuns (\code{firmware/common/})}} \\
\code{pcb} & mapa de pinos da PCB única, conferido na compilação \\
\code{protocol} & o contrato em C: tópicos MQTT, identificadores CAN, mensagens ESP-NOW \\
\code{rtos\_util} & prioridades nomeadas, núcleos e tamanhos de pilha das tasks \\
\code{sys\_events} & \textit{event group} com o estado do sistema: rede, hora, MQTT, CAN \\
\code{net} & Wi-Fi STA, reconexão com espera crescente, SNTP \\
\code{mqtt\_link} & cliente MQTT 3.1.1 com \textit{Last Will}, fila de saída e novas tentativas \\
\code{can} & transporte TWAI: fila de transmissão, recuperação de \textit{bus-off}, filtro e diretório de IDs \\
\code{indicators} & LED WS2812 numa task própria \\
\midrule
\multicolumn{2}{@{}l}{\textit{Hub (\code{firmware/room/})}} \\
\code{sensors} & eventos das entradas com filtro de \textit{glitch} e \textit{debounce}; drivers do PIR, das portas e do LM35 \\
\code{rules} & motor de regras em C puro, testado no computador \\
\code{hub\_config} & configuração em JSON recebida por MQTT, validada e gravada na NVS \\
\code{espnow\_link} & recepção ESP-NOW: validação, descarte de reenvios, remetente pareado \\
\midrule
\multicolumn{2}{@{}l}{\textit{Vestível (\code{firmware/gadget/})}} \\
\code{mpu6050} & driver do sensor: configuração, FIFO, detecção de sensor perdido \\
\code{fall\_pipeline} & pipeline de queda em C puro, gerado e testado contra o Python \\
\code{signal\_proc} & buffer circular, momentos com ESP-DSP e nível de atividade \\
\code{espnow\_tx} & envio ESP-NOW: varredura de canal, confirmação, reenvio \\
\code{battery} & leitura de um divisor no ADC (para uma placa futura com LiPo) \\
\bottomrule
\end{tabular}
\end{table}

A configuração de cada placa (rede, broker, chaves, MAC do par, opcionais) é feita no
\code{menuconfig} e fica no \code{sdkconfig}, que não vai para o repositório: senhas e chaves
nunca entram no git. O que vale para todos (watchdog, partições) fica num
\code{sdkconfig.defaults} versionado.

\subsection{Tasks, núcleos e recursos compartilhados}\label{sec:tasks}

As tasks seguem o padrão das aulas de RTOS~\cite{slides}: prioridades nomeadas, nunca acima de 20 (a task
do Wi-Fi roda em 23), e cada task fixada num núcleo. O núcleo 0 fica com o rádio e tudo o
que fala com ele; o núcleo 1, com o tempo real (Tabela~\ref{tab:tasks}).

\begin{table}[H]
\centering
\caption{Tasks das duas placas.}
\label{tab:tasks}
\small
\begin{tabular}{@{}llll>{\raggedright\arraybackslash}p{6.6cm}@{}}
\toprule
Placa & Task & Núcleo & Prioridade & Função \\
\midrule
hub & \code{sensors\_task} & 1 & ALTA (15) & estabiliza e lê as portas e o PIR \\
hub & \code{lm35\_task} & 1 & BAIXA (5) & temperatura a cada 30\,s \\
hub & \code{hub\_task} & 0 & MÉDIA (10) & regras, publicação, LED \\
hub & \code{espnow\_task} & 0 & MÉDIA (10) & valida e encaminha as mensagens do vestível \\
as duas & \code{can\_rx}, \code{can\_tx} & 0 & MÉDIA (10) & recepção e transmissão CAN \\
as duas & \code{indicators} & 0 & MÍNIMA (2) & LED \\
vestível & \code{sampling} & 1 & ALTA (15) & FIFO do MPU-6050 e pipeline de queda \\
vestível & \code{comm} & 0 & MÉDIA (10) & ESP-NOW, botão, heartbeat, opcionais \\
\bottomrule
\end{tabular}
\end{table}

Os recursos compartilhados seguem regras simples, que evitam corridas e travamentos:
\begin{itemize}
  \item \textbf{Filas, sempre por valor.} Quem produz um evento nunca toca o estado de quem
        consome. Interrupções e \textit{callbacks} de rádio só enfileiram ou notificam.
  \item \textbf{Event group} para o estado do sistema (rede, hora, MQTT, CAN), com um único
        escritor por bit.
  \item \textbf{Notificação de task} onde um aviso não pode se perder: o \textit{callback} de
        estado do TWAI só acorda a \code{can\_tx}, que lê o estado ela mesma.
  \item \textbf{Mutex} só para somar um lote de atividade no vestível, e seção crítica
        (\code{portMUX}) só para poucas instruções.
  \item Nenhuma task segura duas travas ao mesmo tempo, e as tasks vigiadas pelo
        \textit{watchdog} de 30\,s só esperam com tempo limite.
\end{itemize}

\subsection{Hub do quarto}

\paragraph{Fluxo.} Tudo o que acontece no quarto vira um evento na fila do hub, e uma única
task, a \code{hub\_task}, roda as regras e publica (Figura~\ref{fig:tasks_hub}). Por isso o
motor de regras não precisa de travas. As mudanças de porta e de presença e cada temperatura
saem por MQTT e pelo CAN; as mensagens do vestível chegam pela \code{espnow\_task}, que as
valida antes (versão, tipo, tamanho, remetente pareado e reenvios).

\begin{figure}[H]
\centering
\resizebox{\linewidth}{!}{% Hub: fontes de eventos -> fila do hub -> hub_task (única que roda as regras e publica).
\begin{tikzpicture}[font=\scriptsize, node distance=4mm]
  \tikzset{task/.style={bloco=#1, minimum width=2.5cm, minimum height=0.8cm},
           isr/.style={draw=black!50, fill=black!4, rounded corners=1pt, align=center, minimum width=2.1cm, minimum height=0.7cm},
           fila/.style={draw=black!60, fill=white, minimum width=0.9cm, minimum height=0.42cm, font=\tiny,
                        path picture={\foreach \k in {1,2,3} \draw[black!40] ($(path picture bounding box.west)!\k/4!(path picture bounding box.east)$) ++(0,-0.2) -- ++(0,0.4);}}}
  % fontes
  \node[isr] (isr) at (0,3.0) {ISR do GPIO\\(portas, PIR)};
  \node[task=corhub] (sens) at (4.0,3.0) {\textbf{sensors\_task}\\núcleo 1, ALTA};
  \node[task=corhub] (lm) at (4.0,1.8) {\textbf{lm35\_task}\\núcleo 1, BAIXA};
  \node[task=corambas] (esp) at (4.0,0.6) {\textbf{espnow\_task}\\núcleo 0, MÉDIA};
  \node[isr] (espcb) at (-0.2,0.6) {callback de\\recepção ESP-NOW};
  \node[isr] (mq) at (4.0,-0.6) {MQTT\\\code{config/desejada}};
  \node[fila] (q1) at (1.95,3.0) {};
  \node[fila] (q0) at (1.95,0.6) {};
  \draw[seta] (isr) -- (q1); \draw[seta] (q1) -- (sens);
  \draw[seta] (espcb) -- (q0); \draw[seta] (q0) -- (esp);
  % fila do hub e hub_task
  \node[fila, minimum width=1.4cm, minimum height=0.55cm] (qh) at (6.9,1.2) {};
  \node[font=\tiny, above=0pt of qh] {fila do hub};
  \node[task=corhub, minimum width=2.8cm, minimum height=1.2cm] (hub) at (9.5,1.2) {\textbf{hub\_task}\\núcleo 0, MÉDIA\\regras + publicação};
  \foreach \n in {sens,lm,esp,mq} {\draw[seta] (\n.east) -- ++(0.3,0) |- (qh.west);}
  \draw[seta] (qh) -- (hub);
  % saídas
  \node[bloco=corcentral, minimum width=2.2cm] (mqtt) at (13.0,2.6) {MQTT\\(estados, alertas)};
  \node[bloco=cinza, minimum width=2.2cm] (led) at (13.0,1.2) {LED\\(\code{indicators})};
  \node[fila] (qc) at (12.0,-0.2) {};
  \node[task=corambas, minimum width=2.1cm] (ctx) at (13.8,-0.2) {\textbf{can\_tx}\\TWAI};
  \draw[seta] (hub.east) -- ++(0.3,0) |- (mqtt.west);
  \draw[seta] (hub.east) -- (led.west);
  \draw[seta] (hub.east) -- ++(0.3,0) |- (qc.west);
  \draw[seta] (qc) -- (ctx);
  % recepção CAN
  \node[isr] (twai) at (6.9,-1.6) {ISR do TWAI};
  \node[fila] (qr) at (8.9,-1.6) {};
  \node[task=corambas, minimum width=2.1cm] (crx) at (10.9,-1.6) {\textbf{can\_rx}\\log};
  \draw[seta] (twai) -- (qr); \draw[seta] (qr) -- (crx);
  % legenda
  \node[fila, anchor=west] (lf) at (0,-1.6) {};
  \node[anchor=west, font=\tiny] at (lf.east) {fila (por valor)};
\end{tikzpicture}
}
\caption{Tasks e filas do hub. A \code{hub\_task} é a única que chama as regras e publica.}
\label{fig:tasks_hub}
\end{figure}

\paragraph{Regras.}\label{sec:regras} O motor de regras (Tabela~\ref{tab:regras}) é C puro,
sem dependência do ESP-IDF, e roda também no computador, onde é testado. Cada alerta é um
\textbf{episódio}: publicado uma vez como \code{ativo} e encerrado como \code{resolvido} ou
\code{cancelado}, sem repetição (Figura~\ref{fig:alerta_estados}). Sem hora válida, as regras
com janela de horário ficam suspensas, e o \code{status} do hub informa
\code{"hora":false}. A inatividade só vale de dia porque, à noite, o residente dormindo não
aciona o PIR.

\begin{table}[H]
\centering
\caption{Regras de alerta do hub (valores padrão, ajustáveis por MQTT).}
\label{tab:regras}
\small
\begin{tabular}{@{}l>{\raggedright\arraybackslash}p{4.9cm}>{\raggedright\arraybackslash}p{3.8cm}>{\raggedright\arraybackslash}p{1.9cm}@{}}
\toprule
Tipo & Dispara quando & Resolve quando & Padrão \\
\midrule
\code{queda} & chega uma \code{FALL} do vestível & ``estou bem'' (\code{cancelado}) & --- \\
\code{saida\_noturna} & a porta de saída abre à noite e passam N\,min sem presença & a presença volta & 22h--6h, 15\,min \\
\code{banheiro\_prolongado} & a porta do banheiro fecha e passam N\,min sem presença e sem a saída abrir & a porta do banheiro abre & 40\,min \\
\code{inatividade} & de dia, N\,min sem presença e sem a saída abrir & a presença volta & 120\,min \\
\code{temperatura} & fora da faixa por 10\,min seguidos & dentro da faixa (histerese de 0,5\,°C) por 10\,min & 18--28\,°C \\
\code{vestivel\_offline} & 90\,s sem mensagens do vestível & a próxima mensagem & --- \\
\code{bateria\_baixa} & bateria do vestível abaixo do limite & acima de 3,6\,V & 3,5\,V \\
\bottomrule
\end{tabular}
\end{table}

\begin{figure}[H]
\centering
\resizebox{0.9\linewidth}{!}{% Ciclo de vida de um episódio de alerta no hub.
\begin{tikzpicture}[font=\scriptsize, >=Stealth]
  \tikzset{estado/.style={draw=#1, fill=#1!10, thick, rounded corners=8pt, minimum width=2.2cm, minimum height=0.9cm, align=center}}
  \node[estado=cinza] (ok) at (0,0) {normal\\(sem episódio)};
  \node[estado=corvest] (at) at (5.0,0) {\textbf{ativo}\\LED vermelho};
  \node[estado=corcentral] (res) at (10.6,1.4) {\textbf{resolvido}};
  \node[estado=corcentral] (can) at (10.6,-1.4) {\textbf{cancelado}};
  \draw[->, thick] (ok) -- node[rotulo, above=2pt] {condição da regra} node[rotulo, below=2pt] {publica \code{ativo}} (at);
  \draw[->, thick] (at.east) ++(0,0.2) -- node[rotulo, above=3pt, sloped] {a condição acabou} (res.west);
  \draw[->, thick] (at.east) ++(0,-0.2) -- node[rotulo, below=3pt, sloped] {``estou bem'' (só queda)} (can.west);
  \draw[->, thick, black!50] (res.north) to[out=150, in=40] node[rotulo, above] {publica \code{resolvido} e volta ao normal} (ok.north);
  \draw[->, thick, black!50] (can.south) to[out=210, in=-40] node[rotulo, below] {publica \code{cancelado} e volta ao normal} (ok.south);
\end{tikzpicture}
}
\caption{Ciclo de vida de um episódio de alerta.}
\label{fig:alerta_estados}
\end{figure}

\paragraph{Configuração sem recompilar.} Os parâmetros das regras chegam pelo tópico
\code{config/desejada}, em JSON. O hub valida tudo, grava na NVS e responde em
\code{config/aplicada}; uma configuração inválida é recusada inteira, e a anterior continua
valendo. No boot, o hub usa a NVS ou, se ela estiver vazia, os valores padrão:
\begin{lstlisting}[language=json]
{"v":1, "noite":{"inicio":"22:00","fim":"06:00"},
 "portas":{"quarto":"saida","banheiro":"banheiro"},
 "saida_noturna_min":15, "banheiro_min":40, "inatividade_min":120,
 "temperatura":{"min":18,"max":28,"histerese":0.5,"persistencia_min":10},
 "vestivel_offline_s":90, "bateria_baixa_mV":3500}
\end{lstlisting}

\subsection{Vestível}

\paragraph{Aquisição.} A task \code{sampling} lê a FIFO do MPU-6050 a cada 50\,ms e passa
cada amostra pelo pipeline de queda (Figura~\ref{fig:tasks_vestivel};
Seção~\ref{sec:queda}). Se a FIFO transbordar ou a leitura falhar no meio, a FIFO é zerada e
o pipeline recomeça: uma janela nunca junta trechos dos dois lados de uma lacuna. Dez leituras
seguidas sem amostra (0,5\,s) significam que o sensor parou (um MPU que reinicia volta
dormindo); o vestível passa então a procurar o sensor a cada 5\,s.

\begin{figure}[H]
\centering
\resizebox{\linewidth}{!}{% Vestível: sampling (núcleo 1) roda o pipeline; comm (núcleo 0) cuida do rádio.
\begin{tikzpicture}[font=\scriptsize]
  \tikzset{etapa/.style={draw=corvest, fill=corvest!8, rounded corners=2pt, align=center, minimum height=0.8cm,
                         minimum width=1.55cm, inner sep=2pt},
           fila/.style={draw=black!60, fill=white, minimum width=0.9cm, minimum height=0.42cm, font=\tiny,
                        path picture={\foreach \k in {1,2,3} \draw[black!40] ($(path picture bounding box.west)!\k/4!(path picture bounding box.east)$) ++(0,-0.2) -- ++(0,0.4);}}}
  \node[draw=black!50, fill=black!4, rounded corners=1pt, align=center, minimum height=0.8cm] (mpu) at (-0.7,0)
       {MPU-6050\\FIFO (200\,Hz)};
  \node[etapa] (f) at (2.4,0) {passa-baixa\\20\,Hz (6 eixos)};
  \node[etapa] (g) at (4.4,0) {gatilho\\$\amag \ge 2$\,g};
  \node[etapa] (j) at (6.4,0) {janela\\$-$1,0 a +0,8\,s};
  \node[etapa] (a) at (8.4,0) {23\\atributos};
  \node[etapa] (c) at (10.4,0) {Random\\Forest};
  \node[etapa] (d) at (12.4,0) {$P \ge$ limiar?\\cooldown 30\,s};
  \draw[seta] (mpu) -- node[rotulo, below=9pt] {I²C, a cada 50\,ms} (f);
  \foreach \u/\v in {f/g, g/j, j/a, a/c, c/d} {\draw[seta] (\u) -- (\v);}
  \begin{scope}[on background layer]
    \node[draw=corvest, dashed, rounded corners, fit=(f)(d), inner sep=3mm,
          label={[font=\scriptsize\bfseries, text=corvest]above:\textbf{sampling}: núcleo 1, prioridade ALTA}] (samp) {};
  \end{scope}
  % comm
  \node[fila] (q) at (12.4,-1.9) {};
  \node[bloco=corvest, minimum width=3.0cm] (comm) at (9.4,-1.9) {\textbf{comm}: núcleo 0, MÉDIA\\ESP-NOW, heartbeat, botão};
  \draw[seta] (d) -- node[rotulo, right] {queda} (q);
  \draw[seta] (q) -- (comm);
  \node[draw=black!50, fill=black!4, rounded corners=1pt, align=center] (bt) at (5.6,-1.9) {botão\\(ISR)};
  \draw[seta] (bt) -- (comm);
  \node[bloco=corhub, minimum width=2.2cm] (hub) at (9.4,-3.5) {hub\\(ESP-NOW cifrado)};
  \draw[sem fio=corvest] (comm) -- node[rotulo, right] {\code{FALL}, \code{FALL\_CANCEL}, \code{HEARTBEAT}} (hub);
  \node[draw=black!40, dashed, rounded corners, align=center, font=\tiny] (op) at (5.6,-3.5)
       {opcionais: MQTT direto\\e CAN de bancada};
  \draw[seta=black!40, dashed] (comm.west) ++(0,-0.3) -| (op.north);
\end{tikzpicture}
}
\caption{O vestível: a task \code{sampling} roda o pipeline amostra a amostra no núcleo 1, e a
\code{comm} cuida do rádio no núcleo 0.}
\label{fig:tasks_vestivel}
\end{figure}

\paragraph{Comunicação e interação.} A \code{comm} manda a \code{FALL} na hora em que o
pipeline decide queda, e o LED fica em alerta durante a janela de cancelamento de 15\,s; se o
botão for apertado nesse tempo, sai uma \code{FALL\_CANCEL}. A cada 30\,s sai um
\code{HEARTBEAT} com a bateria e o nível de atividade, o valor RMS de $\amag - 1$\,g desde o
heartbeat anterior. Os detalhes do enlace estão na Seção~\ref{sec:espnow}.

\paragraph{Modos.}
\begin{itemize}
  \item \textbf{Normal}: ESP-NOW até o hub.
  \item \textbf{Coleta}: cada amostra bruta e cada janela classificada (com os 23 atributos e
        a probabilidade) vão para a serial. É assim que vai ser formado o conjunto de dados
        no pulso (Seção~\ref{sec:queda}).
  \item \textbf{Sem hub}: com o MAC do hub vazio na configuração, o ESP-NOW não é ligado e as
        quedas vão só para o log. É o modo de bancada e de demonstração com uma placa só,
        ligada ao computador pela USB (Seção~\ref{sec:demo}).
  \item \textbf{Opcionais}, desligados por padrão: MQTT direto ao broker, como reserva se o
        hub não confirmar uma queda, e CAN de bancada.
\end{itemize}

\subsection{Robustez}\label{sec:ausente}

Uma regra vale para as duas placas: \textbf{uma peça ausente gera um aviso, não um erro}.
Sensor, CAN, Wi-Fi ou broker que faltam dão um único aviso no log quando a falta é percebida e
uma informação quando voltam; as novas tentativas ficam em silêncio, com espera crescente até
60\,s. Nada opcional pode derrubar a placa:
\begin{itemize}
  \item \textbf{Vestível sem MPU-6050}: segue rodando, com o LED vermelho lento, e \textbf{não
        manda heartbeat}, para não parecer saudável sem detectar quedas; o hub o declara
        offline em 90\,s.
  \item \textbf{Hub com sensor não instalado}: um PIR solto lê ``ninguém'' e uma porta solta
        lê ``aberta'', e nada no sinal os distingue de um caso real. Por isso os sensores
        instalados são declarados no \code{menuconfig}. Sem PIR, a inatividade e a saída
        noturna ficam suspensas; uma porta ausente não é lida.
  \item \textbf{Rede}: o Wi-Fi reconecta com espera de 1 a 60\,s; o cliente MQTT só inicia
        com a rede de pé e é reiniciado se o início falhar.
  \item \textbf{Queda que não chegou}: a \code{FALL} é guardada e reenviada igual até ser
        entregue (Seção~\ref{sec:espnow}).
  \item \textbf{Configuração errada não compila}: pino repetido ou proibido, chave do ESP-NOW
        com tamanho diferente de 16 caracteres ou MAC mal escrito fazem o build falhar, com a
        mensagem do problema.
\end{itemize}

\subsection{Testes no computador}

A lógica que não depende de hardware é C puro e roda também no computador, compilada com o
gcc. Cada componente tem uma bateria de testes (Tabela~\ref{tab:testes}); os do pipeline de
queda comparam o C, janela a janela, com o Python que treinou o modelo.

\begin{table}[H]
\centering
\caption{Testes automáticos que rodam no computador (sem placa).}
\label{tab:testes}
\small
\begin{tabular}{@{}l>{\raggedright\arraybackslash}p{8.8cm}r@{}}
\toprule
Componente & O que verifica & Verificações \\
\midrule
\code{rules} (hub) & cada regra dispara e resolve como especificado, episódios sem repetição, queda depois de o vestível reiniciar, sensores ausentes & 62 \\
\code{hub\_config} (hub) & validação do JSON da configuração, recusa de valores inválidos & 60 \\
\code{fall\_pipeline} (vestível) & filtro, gatilho e, em 26 janelas de 8 trechos da SisFall, os mesmos picos, atributos e probabilidades do Python & 885 \\
\code{signal\_proc} (vestível) & buffer circular, atributos, nível de atividade & 25 \\
\midrule
notificador (central) & formatação, reenvio ao Telegram, dados inesperados & 19 testes \\
\code{demo\_vestivel.py} & leitura da serial, reconexão, linhas estragadas & 15 testes \\
\code{modelo/} (Python) & blocos do pipeline contra os exemplos das aulas, leitura dos dados & 30 testes \\
\bottomrule
\end{tabular}
\end{table}

\subsection{Uso de memória}

A Tabela~\ref{tab:memoria} mostra o tamanho de cada firmware, medido pelo \code{idf.py size}.
A partição de aplicação tem 1,5\,MB. A maior parte da imagem é a pilha de Wi-Fi, TCP/IP e TLS;
a floresta do modelo de queda ocupa \rfKB\,kB.

\begin{table}[H]
\centering
\caption{Uso de memória (ESP-IDF 6.0.1, \code{idf.py size}).}
\label{tab:memoria}
\small
\begin{tabular}{@{}lrrr@{}}
\toprule
Firmware & Imagem (bytes) & Partição livre & RAM interna usada \\
\midrule
Hub & 992\,672 & 35\,\% & 112\,135 B (32,8\,\%) \\
Vestível com ESP-NOW & 960\,688 & 37\,\% & 130\,615 B (38,2\,\%) \\
Vestível sem hub (bancada) & 383\,184 & 75\,\% & 86\,203 B (25,2\,\%) \\
\bottomrule
\end{tabular}
\end{table}
